\ssscp{Medial *CC clusters in reduplicated monosyllables}{13-02-01-01}
The most common pattern for reduplicated monosyllables with a
medial *CC cluster in the region is like the \ili{Kambera} data noted
by Blust, retaining the second consonant in the cluster (the onset of the
second syllable) and losing the first (the coda of the first syllable): CVC\sub{1}C\sub{2}VC → CVC\sub{2}V(C).
The \ili{Buru} data in example \qf{13-09} are updated
and expanded from \citet{Grimes1991a}.

\noindent{\wg\st{3pt}\begin{tabularx}
{\linewidth}{>{\hspace{-\tabcolsep}}l<{\hspace{-\tabcolsep}}
>{\hspace{-\tabcolsep}\enspace}lll
>{\raggedright\arraybackslash}X}
\lsptoprule
%\tbx{13-09}	&	PMP gloss	&	*-C1C2-	&	\ili{Buru} -C2-	&	\ili{Buru} Gloss	\\	\midrule
%\endfirsthead
\tbx{13-09}	&	PMP gloss	&	*-C1C2-	&	\ili{Buru} -C2-	&	\ili{Buru} Gloss	\\	\midrule
%\endhead
	&	knock, pound, beat	&	*tuk\tbr{t}uk	&	to\tbr{t}o	&	pound, beat, thrust downward	\\
	&		&		&	ka to\tbr{t}o-l	&	mallet, hammer	\\
	&	chop to pieces	&	*tək\tbr{t}ək	&	te\tbr{t}e	&	hack, cut up, chop up, hew	\\
	&	sip, suck	&	*səp\tbr{s}əp	&	se\tbr{s}e	&	sip, nibble, snack with drink	\\
	&	stuff, cram in	&	*sək\tbr{s}ək	&	se\tbr{s}e	&	stuff in, pack down, tamp down	\\
	&	hack, chop to pieces	&	*sak\tbr{s}ak	&	sa\tbr{s}a	&	slash into two sides, hack new path	\\
	&	slip down, slip off	&	*lus\tbr{l}us	&	lu\tbr{l}u	&	slide on bottom (e.g. down muddy slope)	\\
	&	rim, edge, border	&	*biR\tbr{b}iR	&	fi\tbr{f}i-n	&	edge, side, shore, bank (of river)	\\
	&	rasp, file	&	*kiR\tbr{k}iR	&	ki\tbr{k}i	&	scrape, grate (generic action)	\\
	&	grasp in fist	&	*gəm\tbr{g}əm	&	ge\tbr{g}e-n	&	armpit (cf. \tsc{PAnD} *gəlgəl notch)	\\	\hhline{~}
	&		&		&	es-ge\tbr{g}e	&	carry under arm, carry in armpit	\\
	&	deep resounding sound	&	*guŋ\tbr{g}uŋ	&	go\tbr{g}o	&	thunder rumbling (v.)	\\ 	\hhline{~}
	&		&		&	go\tbr{g}o-n	&	thunder (n)	\\
	&	fall, dislodge	&	*tak\tbr{t}ak	&	ta\tbr{t}a-k	&	drop s.t., hand down, leave s.t. behind (cf. \tsc{PPh} (Zorc) *pa-tak-tak drop)	\\
	&	amputated/&	**puD\tbr{p}uD/	&	po\tbr{p}o-k	&	amputate, cut off end, lop	\\ 	\hhline{~}
	&	stump, remnant&	**puŋ\tbr{p}uŋ	&		&	 off\\
	&	box, give a blow	&	**pəg\tbr{p}əg 	&	pe\tbr{p}e	&	tap (drum), touch, restrain momentarily	\\
\lspbottomrule
\end{tabularx}\mbox{}\\}

However, there are also cases of retaining C\sub{1}
and losing C\sub{2}, even in some of the same languages that
also lose C\sub{1} and retain C\sub{2}, as with the \ili{Buru} data in example \qf{13-10},
and sporadic other cases such as \ili{Dhao} *tu\tbr{k}tuk > \it{tu\tbr{k}u} ‘pound,
beat, hit; work metal, smithing’.

\noindent{\wg\st{3pt}\begin{tabularx}
{\linewidth}{>{\hspace{-\tabcolsep}}l<{\hspace{-\tabcolsep}}
>{\hspace{-\tabcolsep}\enspace}lll
>{\raggedright\arraybackslash}X}
\lsptoprule
%\tbx{13-10}	&	Gloss	&	*-C1C2-	&	\ili{Buru} --C1-	&	Gloss	\\	\midrule
%\endfirsthead
\tbx{13-10}	&	Gloss	&	*-C1C2-	&	\ili{Buru} -C1-	&	Gloss	\\	\midrule
%\endhead
	&	rasp, file	&	*ki\tbr{R}kiR	&	ki\tbr{h}i	&	scrape, scratch (generic action); regular *R > \it{h}	\\
	&	shave, scrape off	&	*ki\tbr{s}kis	&	ki\tbr{s}i	&	stroke fingertip, flick, touch, nudge{\footnotemark} 	\\
	&	trim, prune, clear path	&	**ba\tbr{s}bas	&	fa\tbr{s}a	&	cut with sawing motion, trim; decide	\\
\lspbottomrule
\end{tabularx}\mbox{}\\}
\footnotetext{\citet{BlustTrussel2020} assign \ili{Buru} \it{kisi}
to PCMP *kisi ‘pick out with fingers’,
based on \ili{Ngadha} \it{kisi} ‘peel or pare off,
strip or draw off with the fingers’; \ili{Fordata} \it{kisi} ‘pinch,
hold between thumb and forefinger’; \ili{Buru} \it{kisi-h} ‘remove
(as eye-matter with the finger),
poke out, pry out, de-bone’.
The first bit of the \ili{Buru} gloss ‘remove eye-matter with finger’
is consistent with the primary meaning of the term involving
stroking the index finger across something \citep{GrimesGrimes2020},
but Grimes has not encountered anything like the glosses of ‘poke
out, pry out, de-bone’ associated with this fairly frequent word.}

PMP *ki\tbb{R}\tbr{k}iR ‘rasp,
file’ > \ili{Buru} ki\tbb{h}i ‘scrape,
scratch’ and ki\tbr{k}i ‘scrape, grate’ (near synonyms,
often used for the same action with the same object) show there
can be more than one reflex in the same language from the same etymon.

Similar variation is found in whole branches of subgroups.
In the \tsc{Rote-Meto} branch of \tsc{Timor-Babar},
the usual pattern -- like \ili{Kambera}
and \ili{Buru} -- is to lose the first consonant.
Examples are given in \trf{13-06}.
Note that consonants that are normally lost in \tsc{Rote-Meto}
(*R, *q, *h) are not considered here.

\begin{table}\tcp{\tsc{Rote-Meto} CVC\sub{1}C\sub{2}VC > CVC\sub{2}VC}{13-06}
\begin{adjustbox}{width=\textwidth}
	\begin{threeparttable}\st{2pt}
	\begin{tabular}{llllll}\lsptoprule
PMP gloss	&	{\hr}bubble, boil	&	{\hr}weevil	&	{\hr}shake	&	{\hr}shave, scrape off	&	{\hr}gnat, sandfly\\
PMP	&	\hr*buk\tbr{b}uk₁	&	\hr*buk\tbr{b}uk₃	&	\hr*gər\tbr{g}ər	&	\hr*kis\tbr{k}is	&	\hr*ñik\tbr{ñ}ik\\	\midrule
PRM gloss	&	bubble, boil	&	weevil	&	shake	&	brush teeth	&	mosquito\\
PRM	&	**ɓu\tbr{ɓ}u	&	**ka-fu\tbr{f}u-k	&	**na-ŋge\tbr{ŋg}o	&	**na-saki\tbr{k}i	&	**ni\tbr{n}i-k\\	\midrule
\ili{Rikou}	&		&	{\hr\hr}fu\tbr{f}u-ʔ	&	{\hr\hr}na-ke\tbr{k}o	&		&	{\hr\hr}ni\tbr{n}i-ʔ\\
\ili{Bokai}	&	{\hr\hr}bu\tbr{b}u	&	{\hr\hr}fu\tbr{f}u-k	&	{\hr\hr}na-ŋe\tbr{ŋ}o	&	{\hr\hr}na-sa-ki\tbr{k}i	&	{\hr\hr}ni\tbr{n}i-k\\
\ili{Termanu}	&	{\hr\hr}na-sa-bu\tbr{b}u	&	{\hr\hr}fu\tbr{f}u-k	&	{\hr\hr}na-ŋge\tbr{ŋ}o	&	{\hr\hr}na-sa-ki\tbr{k}i, ki\tbr{k}i	&	{\hr\hr}ni\tbr{n}i-k\\
\ili{Tii}	&		&	{\hr\hr}fu\tbr{f}u-k	&	{\hr\hr}na-ŋge\tbr{ŋg}o	&	{\hr\hr}(na-saki\tbr{k}i ?)	&	{\hr\hr}ni\tbr{n}i-k\\
\ili{Dengga}	&	{\hr\hr}bu\tbr{b}u	&	{\hr\hr}fu\tbr{f}u-ʔ	&	{\hr\hr}na-ŋge\tbr{ŋg}o	&		&	{\hr\hr}ni\tbr{n}i-ʔ\\
\ili{Oenale}	&	{\hr\hr}na-sa-ɓu\tbr{ɓ}u	&	{\hr\hr}fu\tbr{f}u-ʔ	&	{\hr\hr}na-ŋge\tbr{ŋg}o	&	{\hr\hr}na-sa-ʔi\tbr{ʔ}i\su{†}	&	{\hr\hr}ni\tbr{n}i-ʔ\\
\ili{Kotos} \ili{Amarasi}	&		&	{\hr\hr}ʔfu\tbr{f}uʔ	&		&	{\hr\hr}na-ski\tbr{k}i	&	\hr\cg{(basi, muut)}\\
\ili{Molo}	&	{\hr\hr}n-bu\tbr{b}u	&	{\hr}‹fu\tbr{f}uk›	&	{\hr}‹na-ke\tbr{k}u›, &	{\hr\hr}na-ski\tbr{k}i	&\\
			&										&											&	{\hr}‹na-ke\tbr{k}o›	&		&	\\
		\lspbottomrule
	\end{tabular}
		\begin{tablenotes}\tnsize
			\item[†] {\ili{Dela} has \it{na-saʔiʔi} is ‘use tobacco to rub one’s teeth after chewing betel nut’.}
		\end{tablenotes}
	\end{threeparttable}
\end{adjustbox}
\end{table}

However,
there are also a number of PMP reduplicated monosyllables in
which the first consonant is retained
and the second consonant is lost,
as shown in \trf{13-07}.
Often in such cases there is a split with the PMP form producing
two reflexes in \tsc{Rote-Meto}; one with the first consonant
retained and one with the second consonant retained.

\begin{table}\tcp{\tsc{Rote-Meto} CVC\sub{1}C\sub{2}VC > CVC\sub{1}VC {\tl} CVC\sub{2}VC}{13-07}
\begin{adjustbox}{width=\textwidth}
	\begin{threeparttable}\st{2pt}
	\begin{tabular}{lll|ll|ll}\lsptoprule
PMP gloss	&	\mc{2}{l|}{stuff, cram in; be crowded}	&	\mc{2}{l|}{knock, pound, beat}	&	\mc{2}{l}{pull at body part; hold\su{Δ}}\\
PMP	&	\mc{2}{l|}{*sə\tbb{k}\tbr{s}ək}	&	\mc{2}{l|}{*tu\tbb{k}\tbr{t}uk}	&	\mc{2}{l}{*bi\tbb{t}bit}\\	\midrule
PRM gloss	&	{\hr}stuff, cram in	&	{\hr}force, stuck	&	{\hr}beat, pound	&	{\hr}beat, knock	&	{\hr}pinch	&	{\hr}carry\su{ǁ}\\
PRM	&	**se\tbr{s}ə	&	**se\tbb{k}e\su{‡}	&	**tu\tbr{t}u	&	**to\tbb{k}o	&	**ɓi\tbb{t}i	&	**fi\tbb{t}i\\	\midrule
\ili{Rikou}	&	{\hr\hr}se\tbr{s}e	&	{\hr\hr}na-se\tbb{k}e	&	{\hr\hr}tu\tbr{t}u	&	{\hr\hr}to\tbb{ʔ}o	&		&	\\
\ili{Bokai}	&	{\hr\hr}se\tbr{s}e	&	{\hr\hr}na-ka-se\tbb{ʔ}e	&	{\hr\hr}tu\tbr{t}u	&	{\hr\hr}to\tbb{k}o	&		&	\\
\ili{Termanu}	&	{\hr\hr}se\tbr{s}e	&	{\hr}(na-ka-)se\tbb{ʔ}e	&	{\hr\hr}tu\tbr{t}u	&	{\hr\hr}to\tbb{k}o	&	{\hr\hr}bi\tbb{ʔ}i	&	{\hr\hr}fi\tbb{ʔ}i\\
Ba{\Q}a	&		&		&		&		&		&	{\hr\hr}fi\tbb{ʔ}i\\
\ili{Tii}	&	{\hr\hr}se\tbr{s}e	&	{\hr\hr}na-ka-se\tbb{ʔ}e	&	{\hr\hr}tu\tbr{t}u	&	{\hr\hr}to\tbb{k}o	&	{\hr\hr}ɓi\tbb{ʔ}i	&	{\hr\hr}fi\tbb{ʔ}i\\
\ili{Dengga}	&	{\hr\hr}se\tbr{s}a	&	{\hr\hr}na-ʔa-se\tbb{ʔ}e	&	{\hr\hr}tu\tbr{t}u	&	{\hr\hr}to\tbb{ʔ}o	&	{\hr\hr}bi\tbb{ʔ}i	&	\\
\ili{Oenale}	&	{\hr\hr}se\tbr{s}a	&	{\hr\hr}na-ʔa-se\tbb{ʔ}e	&	{\hr\hr}tu\tbr{t}u	&	{\hr\hr}to\tbb{ʔ}o	&		&	\\
K. \ili{Amarasi}	&	{\hr\hr}na-ʔsee\tbr{s}\su{†}	&	{\hr\hr}na-ʔse\tbb{k}eʔ	&	{\hr\hr}n-tu\tbr{t}u	&		&	{\hr\hr}kbi\tbb{t}i	&	{\hr\hr}n-fi\tbb{t}i\\
\ili{Molo}	&		&	{\hr}‹n(a)-se\tbb{k}e›	&	{\hr\hr}tu\tbr{t}u	&		&	{\hr}‹bi\tbb{t}i›, kabi\tbb{t}i	&	{\hr\hr}n-fi\tbb{t}i\\
\ili{Kusa-Manea}	&		&		&		&		&	{\hr\hr}bi\tbb{t}i	&	{\hr\hr}fii\tbb{ʔ}\\
		\lspbottomrule
	\end{tabular}
		\begin{tablenotes}\tnsize
			\item[†]	{\ili{Kotos} \ili{Amarasi} \it{na-ʔsees} ‘crowded,
								tight’ has final CV → VC metathesis.
								The unmetathesised form has not yet been attested.
								It could be *[na-ʔsesa] or *[na-ʔsese].}
			\item[‡]	{\ili{Termanu} \it{na-ka-seʔe,
								seʔe} ‘squeezed/stuck into something’,
								\it{seʔe{\tl}seʔe} ‘completely trapped in something; squeezed, very cramped’.
								\ili{Kotos} \ili{Amarasi} \it{na-ʔsekeʔ} ‘force,
								put pressure on, press to do’.}
			\item[Δ]	{\citet{BlustTrussel2020} reconstruct PMP *bitbit₂ ‘pull at
								body part; hold something dangling from the fingers’.
								In addition to these reflexes,
								there is PRM **bibi ‘pinch’ \citep[100]{Edwards2021}.}
			\item[{ǁ}]	{\ili{Termanu} \it{fiʔi} is glossed as ‘pinch with the tip of
								the fingers or fingernails.
								From that: hold fast with the tips of the fingers’ \citep[135]{Jonker1908},
								\ili{Meto} reflexes are ‘carry an item hanging below the waist (typically
								hanging from a rope or strap)’.
								**t > \it{ʔ} is not the usual reflex of PRM **t in the Rote languages or \ili{Kusa-Manea},
								and may be explicable if we reconstruct medial *tb to Proto-Rote-\ili{Meto}.}
		\end{tablenotes}
	\end{threeparttable}
\end{adjustbox}
\end{table}

To the forms in \trf{13-07} we could add PMP *ŋi\tbb{s}ŋis ‘grin,
show the teeth’ > Proto-Rote-\ili{Meto} **ni\tbb{s}i-k ‘teeth’
\citep[284]{Edwards2021}, which shows a semantic shift that is extremely common for this
word in this region.
Though in this case \citet{BlustTrussel2020} reconstruct the
doublet *ŋisi from which Blust
and Trussel derive the \tsc{Rote-Meto} (and many other Wallacean) forms.

To summarise, while it is true that there is reduction of  medial
*CC clusters across languages in the region with the prevalence
of having no intervocalic medial -CC- clusters in their phonotactics,
nevertheless because C\sub{1} is preserved in some languages
and some items, and C\sub{2} is preserved in other languages
for the same item,
both C\sub{1} and C\sub{2} must be reconstructed to account for
the data across the region,
at l\ili{east} in some etyma.
The occasional pattern of C\sub{1}C\sub{2} → C\sub{3} also requires
reconstruction of the full cluster.
\ili{As} with many other issues in the complex historical phonology
in the region, this must be looked at item-by-item,
and so caution should be taken when making broad generalisations
about medial clusters for any language,
etymon, or subgroup in the region.

Although the most common pattern in the region may be the one
observed by Blust (CVC\sub{1}C\sub{2}VC → CVC\sub{2}VC),
the data from \ili{Buru},
\ili{Dhao}, and \tsc{Rote-Meto} show that the other pattern is also
observed whereby CVC\sub{1}C\sub{2}VC → CVC\sub{1}VC.
This minority pattern is detectable in these languages because
of the great volume of lexical data which is available; \citet{GrimesGrimes2020}
has over 7,071 headwords,
and \citet{Edwards2021} has 1,174 Proto-Rote-\ili{Meto} reconstructions.
For most languages of Wallacea,
this quantity of lexical data is not available,
or where it is available,
PMP etyma have not (yet) been thoroughly checked.
Thus, it is unclear to what extent other languages might also
show CVC\sub{1}C\sub{2}VC → CVC\sub{1}VC.
The evidence from \ili{Buru}
and \tsc{Rote-Meto}, combined with the general lack of data from
other languages, cautions against making broad generalisations.
It is true that reduction is widespread in the region,
but the totality of the data still require reconstruction of the full cluster.

Applying the same principles used for post-nasal voicing (\srf{13-01-02})
and glide truncation (\srf{13-01-03}),
we must conclude that the reduction of *CC clusters in reduplicated
monosyllables was a product of diffusion and/or parallel development,
not of shared inheritance,
since languages and words that preserve C\sub{1} are interspersed
with those that preserve C\sub{2},
and also interspersed with those that coalesce C\sub{1}C\sub{2} → C\sub{3}.
\ili{As} yet, we find no patterns involving these medial *CC clusters
that are helpful or relevant for subgrouping at higher levels.
